%=============================================================================!
% 14.8  ISOTROPIC, SEMI-ISOTROPIC AND ANISOTROPIC COUPLING                    !
%=============================================================================!
\section{Isotropic, semi-isotropic and anisotropic coupling}
\label{sec:pr-shape}

Every barostat run so far has scaled the three lengths of the box by one
factor. That suits a liquid, whose pressure tensor is $P$ times the unit
matrix, and fails for a solid that wants to grow along one direction
only. This section takes graphite filling with lithium as its example,
sets out the ways a barostat can let the shape change, and measures what
each does to a layered model solid.

\subsection*{Graphite opens along one axis}

Graphite is a stack of sheets of carbon, each a honeycomb of strong
bonds, held together by a weak attraction between the sheets. When
lithium enters, it sits between the sheets over the centres of the
hexagons, and the sheets, stacked A, B, A, B in graphite, each shifted so
that half its atoms sit over the centres of the hexagons below, line up
A, A around it, each directly over the last; in the fully lithiated
LiC$_6$ the sheets open from \SI{3.34}{\angstrom} apart to
\SI{3.70}{\angstrom}, 10.8\% wider\cite{Hazrati2014}
(\cref{fig:pr-graphite}). A honeycomb has two carbon atoms for each
hexagon, since each of a hexagon's six corners is shared by three
hexagons, so one lithium for six carbons fills every third hexagon. A barostat that scaled the three lengths together could
open the gaps only by stretching the sheets by as much, and would leave
the solid strained.

\begin{figure}[htb]
\centering
\includegraphics{ch14_pressure/graphite}
\caption{Graphite and LiC$_6$, four sheets of carbon each, seen almost
along the sheets: (a) graphite, stacked A, B, A, B, with the sheets
\SI{3.34}{\angstrom} apart; (b) LiC$_6$, stacked A, A, with lithium
(purple) over every third hexagon between the sheets, \SI{3.70}{\angstrom}
apart (the experimental spacings quoted by Hazrati, de Wijs and
Brocks\cite{Hazrati2014}; the in-plane spacing is kept at graphite's
\SI{2.46}{\angstrom}).}
\label{fig:pr-graphite}
\end{figure}

\subsection*{Three couplings}

For a box whose lattice vectors lie along $x$, $y$ and $z$, a barostat can
respond to the pressure tensor in three ways.
\begin{itemize}
\item \emph{Isotropic}: one factor for all three lengths, from the mean
   of the diagonal, $P$.
\item \emph{Semi-isotropic}: one factor for $x$ and $y$ together, from
   $\frac12(\mathsf{P}_{xx} + \mathsf{P}_{yy})$, and another for $z$, from
   $\mathsf{P}_{zz}$. This suits a layered solid or a
   membrane\cite{Bernetti2020}, whose plane keeps its own symmetry while
   its thickness changes.
\item \emph{Anisotropic}: a factor for each length, from its own diagonal
   component.
\end{itemize}
Changing the angles of the cell as well needs every component of
$\vb{h}$ free, as in Parrinello and Rahman's barostat
(\cref{sec:pr-piston}). \texttt{barostats.Berendsen} takes each of the
three couplings, scaling the length along $\alpha$ by $1 - (\kappa\dt/3
\tau_{\mathrm P})(P_0 - P_\alpha)$ with $P_\alpha$ the pressure that length
responds to; its anisotropic form follows ASE's
\texttt{Inhomogeneous\_NPTBerendsen} to \SI{7.9e-8}{\angstrom} over 40
steps.
Stochastic cell rescaling has a semi-isotropic form too, with its own
drift and noise for the logarithm of the area $A$ normal to $z$ and for
the height $L_z$\cite{Bernetti2020},
\begin{align*}
   \dd\ln A &= -\frac{2\kappa}{3\tau_{\mathrm P}}\Bigl(P_0 -
   \frac{\mathsf{P}_{xx} + \mathsf{P}_{yy}}{2}\Bigr)\dd t + \sqrt{\frac{4\kB T
   \kappa\,\dd t}{3V\tau_{\mathrm P}}}\;g_A ,\\
   \dd\ln L_z &= -\frac{\kappa}{3\tau_{\mathrm P}}\bigl(P_0 - \mathsf{P}_{zz}
   \bigr)\dd t + \sqrt{\frac{2\kB T\kappa\,\dd t}{3V\tau_{\mathrm P}}}\;g_z ,
\end{align*}
with $g_A$ and $g_z$ independent normal numbers. Their sum, since $\ln V =
\ln A + \ln L_z$, is \cref{eq:pr-scr}: the drifts add to
$-(\kappa/\tau_{\mathrm P})(P_0 - P)\,\dd t$, and the variances of the
noises, $\frac23$ and $\frac13$ of $2\kB T\kappa\,\dd t/(V\tau_{\mathrm
P})$, add to the whole of it, since the variances of independent numbers
add (\cref{sec:sm-probability}). The area carries two of the three
lengths and the height one, hence the shares $\frac23$ and $\frac13$.

\subsection*{A layered model solid}

Graphite needs a potential that a later chapter will bring; its point
can be made with Lennard-Jones atoms. The model solid of
\cref{fig:pr-layered} is nine flat sheets of 48 argon-like atoms each, the
atoms of a sheet at the corners of equilateral triangles, stacked so that
each sheet's atoms sit over the hollows of the one below. Atoms within a
sheet attract with eight times argon's depth and 0.95 of its $\sigma$,
atoms in different sheets with argon's own, so the sheets are stiff and
the gaps soft: at 40 K the sheets settle \SI{3.151}{\angstrom} apart, and
the perfect solid, its atoms on their sites in that cell, has $C_{11} =
-\partial\mathsf{P}_{xx}/\partial\epsilon_{xx} = \SI{35.81}{\giga\pascal}$
along the sheets against $C_{33} = -\partial\mathsf{P}_{zz}/\partial
\epsilon_{zz} = \SI{5.31}{\giga\pascal}$ across them, 6.7 times softer.
Between each pair of sheets go twelve guest atoms of lithium's mass, one
for every four sheet atoms, each over a hollow of both sheets, in rows
\SI{7.25}{\angstrom} apart along $x$ and \SI{6.28}{\angstrom} along $y$. They
interact with the sheet atoms through a Lennard-Jones potential of twice
argon's depth, and with each other through one of a fifth of argon's
depth and 1.2 times its $\sigma$. Their reach, the $\sigma$ of their
interaction with the sheets, grows from 1.6 to \SI{2.5}{\angstrom} over
\SI{10}{\pico\second}, pushing the sheets apart through a potential that
changes with time, as the wall pushed in of \cref{sec:en-drift} did, and
is then held for \SI{30}{\pico\second} more, under Berendsen's barostat at
$P_0 = 0$ with each coupling in turn.

\begin{figure}[htb]
\centering
\includegraphics{ch14_pressure/layered}
\caption{Coupling a layered solid. The layered model at \SI{40}{\kelvin}
under Berendsen's barostat at $P_0 = 0$, guests between its sheets whose
reach grows over the first \SI{10}{\pico\second} (dotted) and is then
held for \SI{30}{\pico\second}. (a) The change of the height $c$ of the cell (solid) and of its
in-plane width $a$ (dashed) under isotropic, semi-isotropic and
anisotropic coupling. (b) The pressure in the plane, $\frac12(\mathsf{P}_{xx}
+ \mathsf{P}_{yy})$ (hatched), and across it, $\mathsf{P}_{zz}$ (solid), over the
last \SI{10}{\pico\second}.}
\label{fig:pr-layered}
\end{figure}

Over the last \SI{10}{\pico\second}, under semi-isotropic coupling the
sheets have opened by 7.84\% and the plane has shrunk by 0.20\%, and the
mean pressures in the plane and across it are zero. Under anisotropic
coupling the sheets have opened by 7.42\% and the plane has changed by
less than 0.1\%. Both have settled, the opening wandering by about
0.2\% between successive stretches of \SI{5}{\pico\second}. The two differ
because the rows of guests make the two directions of the plane unequal:
semi-isotropic coupling, which scales $x$ and $y$ together, leaves
$\mathsf{P}_{xx}$ and $\mathsf{P}_{yy}$ at $+0.013$ and
\SI{-0.014}{\giga\pascal}, and anisotropic coupling, which lets each length
move on its own, brings both to zero. Under isotropic coupling every
length grows by 0.97\%: the gaps cannot open as far as the guests need
without the sheets stretching equally, and the solid is left pulled along
its sheets, \SI{-0.141}{\giga\pascal} in the plane, and squeezed across
them, \SI{+0.282}{\giga\pascal}, while the mean of the three, the only
number the isotropic barostat reads, is zero. The early dip of $c$ under
semi-isotropic and anisotropic coupling in \cref{fig:pr-layered}(a), by
about 1\%, comes while the guests' reach is short. Each guest is
\SI{2.60}{\angstrom} from its nearest sheet atoms, which lies beyond the
minimum $2^{1/6}\sigma$ of their pair energy for every $\sigma$ up to
\SI{2.3}{\angstrom}, so until then the guests pull the sheets towards
them.

\begin{keyidea}
A solid that wants to change shape needs a barostat that lets it:
semi-isotropic coupling, one factor for the plane and one across it, for
a layered solid like graphite; anisotropic coupling for each length on its
own, when the directions of the plane differ too. Isotropic coupling of
the layered model with guests between its sheets leaves
\SI{-0.141}{\giga\pascal} along the sheets and \SI{+0.282}{\giga\pascal}
across them while its mean pressure reads zero.
\end{keyidea}

\notebook{14}{insert guests between the sheets under each coupling}

\begin{selfcheck}
A graphite particle opens across its sheets by 10.8\% on lithiation, and
suppose it keeps its in-plane size. By what fraction does its volume grow?
\end{selfcheck}
